\documentclass[../../main.tex]{subfiles}
\begin{document}

{\bf [解]}

% TODO: alpha,betaの追加
\begin{figure}[htb]
  \centering
  \begin{tikzpicture}[scale=2.4, >=stealth, every node/.style={font=\small}]

    % ============================================================
    % 1. 定数・厳密パラメータ
    % 円: x^2 + (y-1)^2 = 1, 中心 (0, 1), 半径 R = 1
    % 直線: x = sqrt(3)/3 ≈ 0.5774
    % y - 1 = ±sqrt(1 - 1/3) = ±sqrt(6)/3 ≈ ±0.8165
    % ============================================================
    \def\xLine{0.57735} % sqrt(3)/3
    \def\yA{0.1835}     % 1 - sqrt(6)/3
    \def\yB{1.8165}     % 1 + sqrt(6)/3

    % 円中心から見た交点の中心角: cos(phi) = sqrt(3)/3 => phi ≈ 54.74°
    % 円の中心は (0, 1) なので、x軸方向の角度は ±54.74°
    \def\phiDeg{54.7356}

    \coordinate (O) at (0, 0);
    \coordinate (A) at (\xLine, \yA);
    \coordinate (B) at (\xLine, \yB);

    % ============================================================
    % 2. 領域 D の斜線塗りつぶし (x >= sqrt(3)/3 かつ 円の内部)
    % ============================================================
    \fill[pattern=north east lines, pattern color=blue!75, opacity=0.75]
    (A) arc (-\phiDeg:\phiDeg:1) -- cycle;

    % ============================================================
    % 3. 座標軸
    % ============================================================
    \draw[->, thick, gray!80!black] (-0.3, 0) -- (1.5, 0) node[right] {$x$};
    \draw[->, thick, gray!80!black] (0, -0.25) -- (0, 2.3) node[above] {$y$};
    \node[below left=2pt] at (O) {$O$};

    % ============================================================
    % 4. 円と直線 x = sqrt(3)/3
    % ============================================================
    % 円全体 (破線)
    \draw[dashed, line width=1.0pt, gray!80!black] (0, 1) circle (1);

    % 円の右側円弧 (領域 D の境界: 実線強調)
    \draw[line width=1.4pt, blue!85!black] (A) arc (-\phiDeg:\phiDeg:1);

    % 直線 x = sqrt(3)/3
    \draw[line width=1.1pt, teal!85!black] (\xLine, -0.15) -- (\xLine, 2.15);
    \node[below=2pt, font=\scriptsize, teal!90!black] at (\xLine, 0) {$\frac{\sqrt{3}}{3}$};

    % 弦 AB (境界線強調)
    \draw[line width=1.5pt, teal!90!black] (A) -- (B);

    % ============================================================
    % 5. 動径 OA, OB および 直線 l
    % ============================================================
    \draw[dashed, gray!60] (O) -- (A);
    \draw[dashed, gray!60] (O) -- (B);

    % 直線 l (θ ≈ 48° 付近の例示)
    \def\thetaVal{48}
    \draw[line width=1.2pt, red!85!black] (O) -- ({2*cos(\thetaVal)}, {2*sin(\thetaVal)})
    node[above right, font=\footnotesize\bfseries] {$l: y = x\tan\theta$};

    % 角度 θ の円弧
    \draw[->, thick, red!85!black] (0.35, 0) arc (0:\thetaVal:0.35);
    \node[red!90!black, font=\scriptsize\bfseries] at (0.42, 0.16) {$\theta$};

    % ============================================================
    % 6. 特徴点とラベル
    % ============================================================
    \fill[black] (A) circle (1.2pt) node[right=2pt] {$A\left(\frac{\sqrt{3}}{3}, 1-\frac{\sqrt{6}}{3}\right)$};
    \fill[black] (B) circle (1.2pt) node[right=2pt] {$B\left(\frac{\sqrt{3}}{3}, 1+\frac{\sqrt{6}}{3}\right)$};

    % y 軸への投影破線と目盛り
    \draw[dotted, gray!50] (0, \yA) -- (A);
    \draw[dotted, gray!50] (0, \yB) -- (B);
    \node[left=2pt, font=\scriptsize] at (0, \yA) {$1-\frac{\sqrt{6}}{3}$};
    \node[left=2pt, font=\scriptsize] at (0, \yB) {$1+\frac{\sqrt{6}}{3}$};

    % 領域ラベル D
    \node[blue!90!black, font=\normalsize\bfseries] at (0.78, 1.0) {$D$};

  \end{tikzpicture}
  \caption{領域 $D$（斜線部）と直線 $l$ および弦 $AB$}
  \label{fig:domain_D_and_line_l}
\end{figure}

与えられた図形を
\begin{align}
  D: x^2+(y-1)^2 \le 1 \land x \ge \dfrac{\sqrt{3}}{3} \label{eq:1}
\end{align}
とし，円と$x=\dfrac{\sqrt{3}}{3}$の交点を
\begin{align}
  A\left(\dfrac{\sqrt{3}}{3}, 1-\dfrac{\sqrt{6}}{3}\right), \quad
  B\left(\dfrac{\sqrt{3}}{3}, 1+\dfrac{\sqrt{6}}{3}\right)
\end{align}
とおく．

題意を満たす $l$ は明らかに傾き正だから
\begin{align}
  l : y = \tan\theta \cdot x, \quad \left(0 < \theta < \dfrac{\pi}{2}\right)
\end{align}
とかける．
以下簡単のため $c = \cos\theta, s = \sin\theta, t = \tan\theta$ と表す．
すると $x$ 軸と $OA, OB$ とのなす角 $\alpha, \beta$ として，
\begin{align}
  \alpha < \theta < \beta \label{eq:2}
\end{align}
であり，$l$ と $x^2+(y-1)^2=1$ との交点のうち$x>0$を満たすもの$P$の$x$座標は連立して
\begin{align}
  (1+t^2)x^2 - 2tx = 0 \quad
  \therefore
  x = \dfrac{2t}{1+t^2} \quad (\because x > 0)
\end{align}
である．$l$と$x=\dfrac{\sqrt{2}}{3}$の交点を$Q$，
$P,Q$から$x$軸に下ろした垂足を$P',Q'$とすれば求める長さ$L$は
\begin{align}
  L & = |PQ| = \dfrac{|P'Q'|}{c} \\
  & = \left( \dfrac{2t}{1+t^2} - \dfrac{\sqrt{3}}{3} \right) / c \\
  & = 2s - \dfrac{\sqrt{3}}{3}\dfrac{1}{c}  \label{eq:3}
\end{align}
である．以下この最大値を求める．

% 領域$D$と直線$l$，L
\begin{figure}[htb]
  \centering
  \begin{tikzpicture}[scale=2.6, >=stealth, every node/.style={font=\small}]

    % ============================================================
    % 1. 定数・厳密パラメータ
    % 例として最大値付近の角 theta0 (cos(theta0) = sqrt(3)/3 => theta0 ≈ 54.74°)
    % ============================================================
    \def\xLine{0.57735} % sqrt(3)/3
    \def\phiDeg{54.7356}

    % theta = 48° で図を描画
    \def\thDeg{48}
    \pgfmathsetmacro{\tVal}{tan(\thDeg)}
    \pgfmathsetmacro{\cVal}{cos(\thDeg)}
    \pgfmathsetmacro{\sVal}{sin(\thDeg)}

    % P の座標: x = 2t/(1+t^2) = 2*s*c, y = x*t = 2*s^2
    \pgfmathsetmacro{\xP}{2.0 * \sVal * \cVal}
    \pgfmathsetmacro{\yP}{2.0 * \sVal * \sVal}

    % Q の座標: x = xLine, y = xLine * t
    \pgfmathsetmacro{\xQ}{\xLine}
    \pgfmathsetmacro{\yQ}{\xLine * \tVal}

    \coordinate (O)  at (0, 0);
    \coordinate (A)  at (\xLine, {1 - sqrt(2.0/3.0)});
    \coordinate (B)  at (\xLine, {1 + sqrt(2.0/3.0)});
    \coordinate (P)  at (\xP, \yP);
    \coordinate (Q)  at (\xQ, \yQ);
    \coordinate (Pp) at (\xP, 0);
    \coordinate (Qp) at (\xQ, 0);

    % ============================================================
    % 2. 領域 D の薄い塗りつぶし
    % ============================================================
    \fill[blue!10, opacity=0.5]
    (A) arc (-\phiDeg:\phiDeg:1) -- cycle;

    % ============================================================
    % 3. 座標軸
    % ============================================================
    \draw[->, thick, gray!80!black] (-0.2, 0) -- (1.4, 0) node[right] {$x$};
    \draw[->, thick, gray!80!black] (0, -0.2) -- (0, 2.2) node[above] {$y$};
    \node[below left=2pt] at (O) {$O$};

    % ============================================================
    % 4. 境界線と直線 l
    % ============================================================
    % 円弧
    \draw[line width=1.2pt, blue!85!black] (A) arc (-\phiDeg:\phiDeg:1);
    \draw[dashed, gray!45] (0, 1) circle (1);

    % 直線 x = sqrt(3)/3
    \draw[line width=1.1pt, teal!85!black] (\xLine, -0.1) -- (\xLine, 2.05);

    % 直線 l 全体
    \draw[line width=1.0pt, gray!40!black] (O) -- ($ 1.25*(P) $) node[above right] {$l$};

    % 角度 θ の円弧
    \draw[->, thick, gray!60!black] (0.35, 0) arc (0:\thDeg:0.35);
    \node[gray!80!black, font=\scriptsize] at (0.42, 0.16) {$\theta$};

    % ============================================================
    % 5. 垂足 P', Q' への垂線と直角直角三角形
    % ============================================================
    \draw[dashed, gray!60] (Q) -- (Qp);
    \draw[dashed, gray!60] (P) -- (Pp);

    % Q からの水平線 (直角三角形の形成)
    \coordinate (Q_horiz) at (\xP, \yQ);
    \draw[dotted, gray!60!black] (Q) -- (Q_horiz);
    \draw[thin, gray!70!black] ($ (Q_horiz) + (-0.08, 0) $) -- ($ (Q_horiz) + (-0.08, 0.08) $) -- ($ (Q_horiz) + (0, 0.08) $);

    % 線分 PQ = L (太い赤線で強調)
    \draw[line width=2.0pt, red!85!black] (Q) -- (P)
    node[midway, above left=2pt, font=\normalsize\bfseries] {$L$};

    % x 軸上の区間 |P'Q'|
    \draw[|<->|, thick, teal!80!black] ($ (Qp) + (0, -0.08) $) -- ($ (Pp) + (0, -0.08) $)
    node[midway, below=2pt, font=\scriptsize\bfseries] {$|P'Q'| = L\cos\theta$};

    % ============================================================
    % 6. 点プロットとラベル
    % ============================================================
    \fill[teal!85!black] (Q) circle (1.3pt) node[above left=1pt] {$Q$};
    \fill[red!85!black]  (P) circle (1.4pt) node[above right=1pt] {$P$};
    \fill[black] (Qp) circle (1.1pt) node[below=2pt, font=\scriptsize] {$Q'\left(\frac{\sqrt{3}}{3}, 0\right)$};
    \fill[black] (Pp) circle (1.1pt) node[below=2pt, font=\scriptsize] {$P'$};

  \end{tikzpicture}
  \caption{領域 $D$ を横切る線分 $PQ = L$ と$x$軸への垂線の足 $P', Q'$ }
  \label{fig:segment_L_projection}
\end{figure}

\cref{eq:3}を$\theta$で微分すると
\begin{align}
  L'(\theta)
  = 2c - \dfrac{\sqrt{3}}{3}\dfrac{s}{c^2}
  = \dfrac{2c^3 - \dfrac{\sqrt{3}}{3}s}{c^2}
\end{align}
である．\cref{eq:2}より$c,s>0$だからこの符号は
\begin{align}
  L'(\theta) \ge 0
  &\iff
  2c^3 \ge \dfrac{\sqrt{3}}{3}s
  \iff
  4c^6 \ge \dfrac{2}{9}s^2 \\
  &\iff
  18c^6 \ge 1-c^2 \\
  &\iff
  (c^2 - \dfrac{1}{3})(18c^4 + 6c^2 + 3) \ge 0 \\
  &\iff
  (c^2 - \dfrac{1}{3})\left\{ 18\left(c^2+\dfrac{1}{6}\right)^2 + \dfrac{5}{2} \right\} \ge 0 \\
  &\iff
  c^2 \ge \dfrac{1}{3} \\
  &\iff
  c \ge \dfrac{\sqrt{3}}{3} \quad (\because c>0)
\end{align}
である．
ここで$\theta_0$ は $0 < \theta_0 < \pi/2 \land \cos\theta_0 = \dfrac{\sqrt{3}}{3}$ をみたすとする．$\alpha,\beta$と$\theta_0$の大小について，
\begin{align}
  \tan\alpha = \dfrac{3\sqrt{2}-\sqrt{14}}{2},
  \tan\beta = \dfrac{3\sqrt{2}+\sqrt{14}}{2},
  \tan\theta_0 = \sqrt{2}
\end{align}
だから，$\tan$の区間内での単調増加性より
\begin{align}
  \alpha < \theta_0 < \beta
\end{align}
である．

以上より$L$の増減表は以下のようになる．

\begin{table}[htb]
  \centering
  \caption{$L$の増減表}
  \begin{tabular}{c|ccccc}
    $\theta$ & $\alpha$ & \dots & $\theta_0$ & \dots & $\beta$ \\ \hline
    $C$ & & & $\dfrac{\sqrt{3}}{3}$ & & \\ \hline
    $L'$ & & $+$ & $0$ & $-$ & \\ \hline
    $L$ & & $\nearrow$ & & $\searrow$ &
  \end{tabular}
\end{table}

したがって，$c = \cos\theta = \dfrac{\sqrt{3}}{3}$ の時に最大値$\max L = \dfrac{\sqrt{6}}{3}$をとる．

\end{document}
